\label{sec:derivabilidad}

\begin{definition}
    \textbf{Derivada parcial.} Sean $A$ un conjunto abierto, $f:A\to \R$ y $\mathbf{x}_0\in \interior{A}$. Se define la \emph{derivada parcial de $f$ con respecto a $x_j$ en $\mathbf{x}_0$}, y se nota por $\partial f/\partial x_j(\mathbf{x}_0)$ o por $f_{x_j}(\mathbf{x}_0)$, a
    \[
        \dif{f}{x_j}(\mathbf{x}_0)=f_{x_j}(\mathbf{x}_0)=\lim_{h\to0}\dfrac{f(\mathbf{x}_0+h\mathbf{e}_j)-f(\mathbf{x}_0)}{h},
    \]
    donde $\mathbf{e}_j$ es el $j$-\'esimo versor can\'onico en $\Rn{n}$ (el vector de $\Rn{n}$ que tiene un 1 en la posición $j$ y ceros en todas las demás), con $j=1,2,\dots,n$. Si esta igualdad no existe, decimos que \emph{$f$ no admite derivada parcial con respecto a $x_j$ en $\mathbf{x}_0$}.
\end{definition}

\begin{definition}
    \textbf{Derivada direccional.} Sean $f:A\subseteq\Rn{n}\to\R$, $\mathbf{x}_0\in \interior{A}$ y $\mathbf{v}=(v_1, v_2, \dots, v_n)$, con $\norm{\mathbf{v}}=1$. Se define la \emph{derivada direccional de $f$ con direcci\'on $\mathbf{v}$ en $\mathbf{x}_0$}, y se nota por $\partial f/\partial \mathbf{v}(\mathbf{x}_0)$ o por $f_{\mathbf{v}}(\mathbf{x}_0)$, a
    \[
        \dif{f}{\mathbf{v}}(\mathbf{x}_0)=f_{\mathbf{v}}(\mathbf{x}_0)=\lim_{h\to0}\frac{f(\mathbf{x}_0+h\mathbf{v})-f(\mathbf{x}_0)}{h}.
    \]
    Si esta igualdad no existe, decimos que \emph{$f$ no admite derivada direccional con direcci\'on $\mathbf{v}$ en $\mathbf{x}_0$}.
\end{definition}

\begin{obs}
    Las derivadas parciales son un caso particular de las derivadas direccionales.
\end{obs}

\begin{obs}
La existencia de las derivadas direccionales no garantiza la continuidad.
\end{obs}

\begin{definition}
    Sea $A \subseteq \mathbb{R}^n$ un conjunto abierto. Una función $f: A \to \mathbb{R}$ es de \emph{clase} $\mathcal{C}^1$ si todas sus derivadas parciales existen y son funciones continuas en $A$.
\end{definition}

\begin{definition}\label{def:ceuno}
    Sea $A\subseteq\Rn{n}$ un conjunto abierto. Una función $f: A \to \mathbb{R}$ es de \emph{clase} $\mathcal{C}^2$ si es de clase $\mathcal{C}^1$ y todas las derivadas parciales de $f$ son de clase $\mathcal{C}^1$.
\end{definition}

\begin{definition}
    Sea $A\subseteq\Rn{n}$ un conjunto abierto. Una función $f: A \to \mathbb{R}$ es de clase $\mathcal{C}^n$ si todas sus derivadas parciales de orden $k$ existen y son continuas para todo $k \leq n$.
\end{definition}

\textbf{Notaci\'on: derivadas de orden superior, derivadas parciales iteradas.} Sea $f:\Rn{2}\to\R$ de clase $\mathcal{C}^2$. Sus \emph{derivadas de segundo orden} se notan de la siguiente manera:
\begin{gather*}
    \frac{\partial}{\partial x}\left(\dif{f}{x}\right)=
    \frac{\partial^2 f}{\partial x^2}=f_{xx},\qquad
    \frac{\partial}{\partial y}\left(\dif{f}{x}\right)=
    \frac{\partial^2 f}{\partial y\partial x}=f_{xy},\\[.3cm]
    \frac{\partial}{\partial y}\left(\dif{f}{y}\right)=
    \frac{\partial^2 f}{\partial y^2}=f_{yy},\qquad
    \frac{\partial}{\partial x}\left(\dif{f}{y}\right)=
    \frac{\partial^2 f}{\partial x\partial y}=f_{yx}.
\end{gather*}
Por supuesto este proceso puede repetirse para derivadas de orden tres y superior. Todas estas se llaman \emph{derivadas parciales iteradas}, mientras que $\partial f/\partial x\partial y$ y $\partial f/\partial y\partial x$ se llaman \emph{derivadas parciales cruzadas}.  

\begin{obs}
    Equivalentemente, la \autoref{def:ceuno} se puede enunciar como: $f$ es de \emph{clase} $\mathcal{C}^2$ si todas sus derivadas parciales de orden dos existen y son continuas.
\end{obs}

\begin{obs}
    Notar que el orden, el cual denota sobre que componentes se aplican las derivadas parciales iteradas, de las componentes en las notaciones de las derivadas de orden superior es invertido. Esto se debe a que $\partial/\partial x$ act\'ua como un operador\footnote{Un operador es una transformaci\'on lineal tal que su dominio e imagen son el mismo espacio vectorial.} sobre $f$, mientras que $f_x$ no.  
\end{obs}

\begin{theorem}\textbf{de Clairaut o de Schwarz, simetr\'ia de las derivadas cruzadas.} 
    
    Sea $A\subseteq\Rn{n}$ un conjunto abierto, $f:A\to\R$ de clase $\mathcal{C}^2$ en $B_\delta(\mathbf{x}_0)$, con $\mathbf{x_0}\in A$. Entonces
    \[
        \dfrac{\partial f}{\partial x_j \partial x_i}(\mathbf{x}_0)=\dfrac{\partial f}{\partial x_i \partial x_j}(\mathbf{x}_0)\quad \text{o}\quad f_{x_i}f_{x_j}(\mathbf{x}_0)=f_{x_j}f_{x_i}(\mathbf{x}_0).
    \]
\end{theorem}

\textcolor{red}{es necesario escribir que A es un conj abierto?}